\chapter{Why the nanoscale is different}
\label{ch:fundamentals}

\paragraph{In this chapter.} What counts as a nanomaterial, why properties change below
$\sim\SI{100}{nm}$ (surface and quantum-size effects), and what we want from a synthesis.
Everything here comes back in the exercises as the \emph{reason} behind an observation
(e.g.\ why a plasmon band shifts, or why small particles are better catalysts).

\section{Definition and general features}
\begin{defbox}[: nanoparticles, clusters, nanocrystals, quantum dots]
Objects of roughly \SIrange{1}{100}{nm} in at least one dimension \slide{8}.
\begin{itemize}
  \item Properties become \emph{discontinuous} as the bulk contracts.
  \item \textbf{Quantum size effects}: electronic, optical, magnetic and mechanical.
  \item \textbf{Large surface effects}; surface modification gives new functionalities.
  \item Difficult to produce by breakdown (top-down) processes; difficult to control shape, size, crystallinity.
  \item Porosity control and dimensionality; not straightforward to characterize; sometimes difficult to handle.
  \item Possible adverse effects on human bodies (nanotoxicity, \slide{16}).
\end{itemize}
\end{defbox}

Historical uses \slide{5--7}: Lycurgus cup (4th c.\ Roman glass with Au/Ag nanoparticles,
different colour in transmitted and reflected light), medieval stained glass (colloidal Au = red,
Ag = yellow), Faraday's gold colloids (1857), Ruska's TEM (1931), Feynman's lecture
``There's plenty of room at the bottom'' (1959), fullerenes (1985), AFM (1986),
IBM logo with 35 Xe atoms (1989), carbon nanotubes (Iijima, 1991).

\section{Quantum size effects}
When a solid shrinks, the continuous bands of the bulk split into discrete levels
\slide{9}: atom $\to$ molecule (HOMO/LUMO) $\to$ cluster $\to$ nanoparticle $\to$ bulk
(semiconductor with a gap, or metal).
\begin{itemize}
  \item \textbf{Band gap opens/increases as size decreases.} This is why CdSe quantum dots
        emit blue when small and red when large \slide{134} (see \cref{ch:uvvis}).
  \item \textbf{Melting point drops} for small particles: Au melts at $\sim\SI{1300}{K}$ in bulk but far lower below
        $\sim\SI{5}{nm}$ (\slide{10}). Small particles are easier to sinter (relevant to catalyst stability, \cref{ch:tem}).
  \item \textbf{Thermal conductivity} of Cu nanoparticles/films decreases for sizes below a few hundred nm \slide{9}.
\end{itemize}

\section{Surface effects}
\begin{keybox}[: surface-to-volume]
Splitting a cube of side \SI{1}{m} into $n^3$ cubes of side $1/n$ multiplies the area by $n$ \slide{12}:
\[
6\times 1^2 = \SI{6}{m^2} \;\to\; 6\times 0.5^2\times 8 = \SI{12}{m^2} \;\to\; 6\times(1/3)^2\times 27 = \SI{18}{m^2}.
\]
In general, $S/V \propto 1/\text{size}$.
\end{keybox}
\begin{itemize}
  \item Fraction of surface atoms in full-shell clusters \slide{13}: 13 atoms $\to$ 92\,\%, 55 $\to$ 76\,\%,
        147 $\to$ 63\,\%, 309 $\to$ 52\,\%, 561 $\to$ 45\,\%, 1415 $\to$ 35\,\%.
        For Fe nanoparticles, surface and bulk atoms are equal in number at $\sim\SI{5}{nm}$.
\item \textbf{Shape effect} \slide{14}: for a cylinder of radius $r$ and height $h$,
        \[ \frac{S}{V} = \frac{2\pi r h + 2\pi r^2}{\pi r^2 h} = 2\left(\frac{1}{r}+\frac{1}{h}\right), \]
        and at fixed volume this is minimal for $h = 2r$ (the slide labels this point ``$h=d$'').
        \checkflag{the plot on slide 14 writes $S/V = 2(r/h + h/r)$ with a minimum at $r/h=1$; that expression is
        not dimensionally a surface-to-volume ratio. The physical conclusion (compact shape $h\approx d$ minimises
        $S/V$) is the same.}
        Rods ($h\gg d$) and discs ($h\ll d$) expose more surface.
\end{itemize}
\begin{tipbox}
``Which particles would you choose for catalysis?'' $\Rightarrow$ the \textbf{smallest, well-dispersed} ones
(highest fraction of surface atoms), provided they are stable (not sintering) and accessible (no thick capping
shell). This argument is used in Exercises~7 and~18.
\end{tipbox}

\section{Applications (overview)}
\slide{17--25}: biomedical (drug delivery, Abraxane, EPR effect = \emph{Enhanced Permeation and Retention} of
nanoparticles in tumours; Au nanoshells for photothermal therapy; Fe$_3$O$_4$ for MRI and magnetic hyperthermia;
Au--antibody conjugates for cancer imaging), antibacterial Ag in textiles and plasters, sunscreens (TiO$_2$, ZnO),
self-cleaning ``lotus'' surfaces, catalysts, energy, electronics.

\section{Objectives of a synthesis}
\label{sec:objectives}
\begin{keybox}[: what a good synthesis should give \slide{27}]
\begin{enumerate}
  \item Monodisperse samples (narrow size distribution).
  \item Same morphology.
  \item Same chemical composition and crystalline structure.
  \item No agglomeration.
  \item If possible: fast, cheap, ``green'', and with \textbf{no critical raw materials}
        (European Critical Raw Materials Act \slide{28}: e.g.\ Pt-group metals, Co, rare earths, \ldots).
\end{enumerate}
\end{keybox}
These five criteria are the yardstick used in \cref{ch:synth-compare} to compare methods.
